% ================================================================
\section{Discussion}\label{sec:discussion}

\textbf{When does the fixed linearization suffice?}
For near-hover operation, TinyMPC's fixed linearization is optimal
by construction.
The flat-parameterized variants become advantageous when the trajectory
induces significant attitude variation---the equilibrium shift
ensures the ADMM solver operates near zero error at every timestep,
improving convergence.

\textbf{Computational hierarchy.}
The five variants form a hierarchy of decreasing online cost
(see also the equivalence chain~\eqref{eq:full_chain}):
\begin{center}
\small
\begin{tabular}{@{}lcccc@{}}
\toprule
Variant & Online ops & Constr. & Conv. & Sched. \\
\midrule
Online Relin & $O(n^3) {+} K O(n^2)$ & Hard & Lin. & $\R^n$ \\
Flat Template & $K \cdot O(n^2)$ & Hard & Lin. & $\R^4$ \\
Barrier + IPM & $K_{\text{I}} \cdot O(n^2)$ & Hard & Quad. & $\R^4$ \\
Barrier & $O(n \cdot m)$ & Soft & --- & $\R^4$ \\
Analytic & $O(n \cdot m)$ & Clip & --- & $\R^4$ \\
\bottomrule
\end{tabular}
\end{center}
The central contribution is the path from cubic-time online
optimization to linear-time deterministic feedback, enabled
by flatness-compressed scheduling over $\R^4$.
The Flat Template has the same per-iteration cost as TinyMPC but
converges in fewer iterations due to better-conditioned linearization.
The Barrier Template (Section~\ref{sec:barrier}) provides soft
constraint satisfaction at the same ${\sim}260$-operation cost as
the Analytic Template, by absorbing log-barrier information
(both $H_\mu$ for inputs and $H_{x,\mu}$ for states)
into the DARE cache.
With optional IPM refinement (2--3 Newton steps),
it achieves hard constraint satisfaction with quadratic convergence---fewer
iterations than ADMM for the same accuracy.
The Analytic Template ($\mu = 0$) trades all constraint handling
for extreme efficiency.
The zero-iteration variants (Barrier and Analytic) achieve
real-linear-time control: deterministic, iteration-free,
$O(n{\cdot}m) \approx 260$ operations per control output,
only possible because flatness compresses the scheduling
space to $\R^4$ (4 sensitivity matrices, 240 floats).

\textbf{What is new vs.\ what is classical.}
The individual mathematical ingredients---differential
flatness~\cite{fliess1995,mellinger2011},
gain scheduling~\cite{rugh2000},
and Riccati sensitivity~\cite{hewer1971,konstantinov1993}---are
well established.
Our contribution is their specific combination in the context
of embedded ADMM-based MPC:
(i)~using flatness to compress the scheduling dimension from
$n$ to~4,
(ii)~caching the \emph{augmented} DARE (with ADMM penalty $\rho$)
rather than the standard LQR DARE,
(iii)~showing the equivalence between the zero-iteration ADMM
forward pass and the flat gain schedule,
and (iv)~absorbing log-barrier constraint penalties into the DARE
cache so that the zero-iteration controller is inherently
constraint-aware (Section~\ref{sec:barrier}).
This combination yields the 240-float, 260-operation
datapath of the Analytic Template---a level of compactness
not achievable by generic gain scheduling or LPV approaches
without the flatness reduction---while the barrier augmentation
adds constraint awareness at no additional online cost.

\textbf{FPGA portability.}
The analytic linearization (Section~\ref{sec:linearize}) and the
Analytic Template (Section~\ref{sec:analytic}) are composed entirely
of elementary operations: trigonometric functions, multiplications,
and additions.
There are no loops with data-dependent iteration counts
and no matrix inversions.
The 240-float storage and ${\sim}260$-operation datapath
are directly realizable as a fixed-function pipeline.

\textbf{Limitations.}
The first-order gain approximation~\eqref{eq:K_interp} degrades
for large $\|\balpha - \balpha_0\|$.
For highly aggressive maneuvers, the Flat Template MPC with a
sufficiently dense cache grid is preferred.
All wall-clock times are from a Python/NumPy reference
implementation and reflect relative computational cost
rather than embedded performance;
C/C++ implementation with Crazyflie hardware experiments
is ongoing work.
